% ============================================================
% Chapter 20 — Statistical Analysis Plan
% Source: PSYCON Engineering Design Document, Vol. 4, Ch. 20
% ============================================================
\section{Statistical Analysis Plan}
\label{sec:statistical-analysis-plan}

Potential analyses include:
\begin{itemize}
    \item Descriptive statistics
    \item Correlation analysis
    \item Classification metrics
    \item Cross-validation
    \item Ablation studies (speech only vs.\ physiology only vs.\ combined)
    \item Error analysis
    \item Confusion matrices
    \item ROC curves (where applicable)
\end{itemize}

\begin{observationbox}
Both overall performance and limitations should be reported.
\end{observationbox}

\subsection{Implemented Evaluation Procedure}
\label{sec:implemented-evaluation-procedure}

The Week 5 evaluator creates one deterministic participant-level train, validation, and test assignment and reuses it for physiology-only, speech-only, and combined comparisons. Preprocessing is fitted inside each training pipeline. Logistic regression and random forest are the current candidates, validation F1 selects the candidate, and the chosen pipeline is refitted on the training and validation participants before one held-out test evaluation.

The output contract includes accuracy, precision, recall, F1, ROC-AUC when both classes are present, confusion matrices, false-positive and false-negative counts, grouped cross-validation, participant-level bootstrap confidence intervals, descriptive statistics, feature correlations, prediction and error rows, modality ablations, data-quality coverage, environment and motion slices, and 10, 30, and 60-second duration summaries. External validation rejects participant overlap and incompatible feature columns.

\begin{observationbox}
These procedures are implemented and covered by unit tests, but no Week 5 metric is reported because no completed marksheet and matching PSYCON recording dataset is available. The target definition, primary metric, \texttt{N/O} handling, exclusions, and duration-stability tolerance must be frozen before the test set is opened.
\end{observationbox}


\subsection{Separate contextual source-score development audit}
The active communication task is ordinal and separate from the Week 5 binary physiology evaluator. The 13 October audit uses 423 source ratings from forty participant-session examples, nine recordings and five known connected groups. The fixed common cohort contains 402 original non-abstained predictions; twenty-one original abstentions remain explicit. Compact grouped development achieved 62.19\% agreement and MAE 0.490, compared with fold training median 64.43\%/0.460 and mode 65.67\%/0.453. New nested pooling exactly matched mode predictions but reduced score-3 recall from 18.45\% to 4.85\%, so it was not promoted. The higher fixed diagnostic is not selected from outer results. No independent final test or 85\% outcome is available.

Training-only preprocessing and inner known-group selection are required. Report both cohorts, per-class precision/recall/F1, confusion matrices, ordinal error, observed-class balanced accuracy and quadratic weighted kappa beside exact agreement. Unresolved people and unreviewed labels prevent participant-independent generalisation and rubric-reliability claims. Full execution records remain private; the aggregate release is \path{docs/report_assets/maximum-defensible-summary.json}.
